% ============================================================
%  main.tex — 学术答辩 Beamer 演示文稿模板
%  ============================================================
%  使用方式：
%    1. 修改下方「元数据」区的标题、作者、导师等信息
%    2. 在各 Section 中替换示例内容为你的正文
%    3. 将图片放入 img/ 目录，用 \includegraphics{img/xxx.png} 引用
%    4. 参考文献放入 refs.bib，用 \cite{key} 引用
%    5. 终端运行 latexmk 即可编译
%  ============================================================

\documentclass[aspectratio=169]{beamer}

% ---------- Style ----------
\usepackage{beamer-style}

% ============================================================
%  元数据 — 修改为你自己的信息
% ============================================================
\title{论文题目：基于 \LaTeX{} Beamer 的\\学术答辩演示文稿模板设计}
\subtitle{硕士研究生学位论文答辩}
\author{张\textbf{三}}
\institute{某某大学 ~ 计算机科学与技术学院}
\date{2026 年 6 月 19 日}

\begin{document}

% ============================================================
%  1. 标题页 (Title Page)
% ============================================================

\begin{frame}[shrink]
  \titlepage

  \vspace{-2em}
  \begin{center}
    \footnotesize
    \textbf{指导教师：}李四 ~ 教授 \\[0.2em]
    \textbf{答辩委员会：}王五 ~ 教授 \quad 赵六 ~ 副教授 \quad 钱七 ~ 研究员
  \end{center}
\end{frame}

% ============================================================
%  2. 目录页 (Table of Contents)
% ============================================================

\begin{frame}{目录}
  \tableofcontents
\end{frame}

% ============================================================
%  3. 正文内容 — 按 Section 组织
% ============================================================

% --- Section: 研究背景 ---
\section{研究背景与意义}

\begin{frame}{研究背景}
  \begin{itemize}
    \item 随着信息技术的快速发展，学术演示文稿已成为知识传播的核心媒介之一。
    \item \LaTeX{} Beamer 因其专业的排版质量和数学公式支持，在理工科领域被广泛使用。
    \item 然而，从零配置 Beamer 主题、字体、配色等工作重复且繁琐，缺乏开箱即用的中文模板。
  \end{itemize}

  \begin{alertblock}{核心问题}
    如何提供一套专业、可复用、支持中文的 Beamer 模板，降低学术答辩演示文稿的制作成本？
  \end{alertblock}
\end{frame}

\begin{frame}{相关工作}
  \begin{columns}[T]
    \begin{column}{0.48\textwidth}
      \textbf{现有方案}
      \begin{itemize}
        \item Overleaf 官方模板库
        \item 各高校自建 \LaTeX{} 模板
        \item 商业软件 (PowerPoint / Keynote)
      \end{itemize}
    \end{column}
    \begin{column}{0.48\textwidth}
      \textbf{局限性}
      \begin{itemize}
        \item 中文支持不完善
        \item 样式与内容耦合，定制困难
        \item 平台锁定，不便于版本管理
      \end{itemize}
    \end{column}
  \end{columns}

  \vspace{1em}
  \begin{exampleblock}{本文方案}
    采用 XeLaTeX + Beamer (metropolis) + PingFang SC 的技术路线，
    实现样式分离、即开即用的中文答辩模板。
  \end{exampleblock}
\end{frame}

% --- Section: 研究方法 ---
\section{研究方法}

\begin{frame}{模板架构设计}
  本模板遵循 \textbf{样式与内容分离} 的设计原则：

  \vspace{0.5em}
  \begin{description}
    \item[\texttt{main.tex}] 用户编辑的主文件，仅包含内容和元数据。
    \item[\texttt{beamer-style.sty}] 样式配置文件，集中管理主题、颜色、字体。
    \item[\texttt{latexmkrc}] 编译配置，一键生成 PDF。
  \end{description}

  \vspace{1em}
  \begin{center}
    \small
    \texttt{用户修改 main.tex → latexmk → 专业排版的 PDF}
  \end{center}
\end{frame}

\begin{frame}{技术选型说明}
  \begin{table}
    \centering
    \begin{tabular}{lll}
      \toprule
      \textbf{组件} & \textbf{选型} & \textbf{理由} \\
      \midrule
      编译器   & XeLaTeX  & 原生 Unicode + 系统字体支持 \\
      Beamer 主题 & metropolis & 简洁现代，适配学术场景 \\
      中文字体 & PingFang SC & macOS 自带，无需额外安装 \\
      参考文献 & biber + biblatex & Unicode 兼容，比 BibTeX 更现代 \\
      \bottomrule
    \end{tabular}
    \caption{技术选型概览}
  \end{table}
\end{frame}

% --- Section: 实验结果 ---
\section{实验结果与分析}

\begin{frame}{编译性能测试}
  \begin{itemize}
    \item 测试环境：macOS + TeX Live 2026，XeLaTeX 编译器
    \item 测试文档：7 页答辩模板（含图表、公式、参考文献）
  \end{itemize}

  \vspace{1em}
  \begin{columns}[T]
    \begin{column}{0.48\textwidth}
      \centering
      \textbf{首次编译} \\
      \vspace{0.3em}
      {\large 2.3 秒}
    \end{column}
    \begin{column}{0.48\textwidth}
      \centering
      \textbf{增量编译} \\
      \vspace{0.3em}
      {\large 0.8 秒}
    \end{column}
  \end{columns}

  \vspace{1em}
  \begin{block}{结论}
    模板编译速度快，latexmk 自动处理多次编译，无需用户干预交叉引用和目录生成。
  \end{block}
\end{frame}

% ============================================================
%  4. 图表页 (Figure Slide)
% ============================================================

\begin{frame}[shrink]{实验结果可视化}
  \begin{figure}
    \centering
    % 占位图：替换为你的实际图片，例如：
    % \includegraphics[width=0.65\textwidth]{img/your-result.png}
    \includegraphics[width=0.55\textwidth]{example-image}
    \caption{实验结果对比图（请替换为实际数据图表）}
    \label{fig:result}
  \end{figure}

  \vspace{-0.8em}
  \begin{itemize}
    \footnotesize
    \item 更多图片示例：\texttt{example-image-a}, \texttt{example-image-b}, \texttt{example-image-golden}
    \item 推荐使用 PDF/PNG 格式矢量图，缩放不失真
  \end{itemize}
\end{frame}

% --- Section: 总结与展望 ---
\section{总结与展望}

\begin{frame}{主要贡献}
  \begin{enumerate}
    \item 设计并实现了一套 \textbf{样式与内容分离} 的 Beamer 中文模板。
    \item 基于 XeLaTeX + metropolis + PingFang SC，确保 \textbf{开箱即用}。
    \item 覆盖学术答辩常见幻灯片类型：标题、目录、正文、图表、参考文献。
    \item 编译命令简化为 \texttt{latexmk}，\textbf{零学习成本}。
  \end{enumerate}
\end{frame}

\begin{frame}{未来工作}
  \begin{itemize}
    \item 支持多套颜色主题切换（亮色 / 暗色 / 学校标准色）。
    \item 提供 Overleaf 兼容的单文件版本。
    \item 集成 minted 宏包实现代码高亮。
    \item 添加 4:3 屏幕比例选项。
  \end{itemize}

  \vspace{1em}
  \begin{center}
    \large 欢迎贡献：\texttt{https://github.com/yourname/beamer-template}
  \end{center}
\end{frame}

% ============================================================
%  5. 致谢页 (Acknowledgments)
% ============================================================

\begin{frame}[c]{}
  \centering
  \Huge 谢谢！\\[0.5em]
  \Large 请各位老师批评指正
\end{frame}

% ============================================================
%  6. 参考文献页 (References)
% ============================================================

\section*{参考文献}

\begin{frame}[allowframebreaks]{参考文献}
  % 列出所有 refs.bib 中的条目（即使未在正文中 \cite）
  \nocite{*}
  \printbibliography
\end{frame}

\end{document}
